\documentclass[12pt,a4paper]{article}

\usepackage[utf8]{inputenc}
\usepackage[T1]{fontenc}
\usepackage[brazil]{babel}
\usepackage{geometry}
\usepackage{graphicx}
\usepackage{float}
\usepackage{hyperref}
\usepackage{xcolor}
\usepackage{listings}
\usepackage{caption}

\geometry{
  left=3cm,
  right=2cm,
  top=3cm,
  bottom=2cm
}

\hypersetup{
  colorlinks=true,
  linkcolor=black,
  urlcolor=blue,
  citecolor=black
}

\lstdefinestyle{sqlstyle}{
  language=SQL,
  basicstyle=\ttfamily\small,
  keywordstyle=\color{blue}\bfseries,
  commentstyle=\color{gray},
  stringstyle=\color{teal},
  numbers=left,
  numberstyle=\tiny\color{gray},
  stepnumber=1,
  numbersep=8pt,
  showstringspaces=false,
  breaklines=true,
  frame=single,
  tabsize=2
}

\newcommand{\inserirDiagrama}[4][0.95\textwidth]{
  \begin{figure}[H]
    \centering
    \IfFileExists{#2}{
      \includegraphics[width=#1]{#2}
    }{
      \fbox{
        \begin{minipage}[c][7cm][c]{0.9\textwidth}
          \centering
          Adicione a imagem \texttt{#2} no Overleaf.
        \end{minipage}
      }
    }
    \caption{#3}
    \label{#4}
  \end{figure}
}

\newcommand{\nomeAluno}{Seu nome}
\newcommand{\matriculaAluno}{Sua matr\'icula}
\newcommand{\cursoAluno}{Seu curso}
\newcommand{\disciplina}{Banco de Dados}
\newcommand{\professor}{Nome do professor}
\newcommand{\instituicao}{Nome da institui\c{c}\~ao}
\newcommand{\tituloTrabalho}{Sistema da Pastoral da Crian\c{c}a para Gerenciamento de Doa\c{c}\~oes}
\newcommand{\localMinimundo}{Pastoral da Crian\c{c}a}
\newcommand{\contatoLocal}{Telefone, e-mail ou endere\c{c}o do local}
\newcommand{\cidadeData}{Cidade -- 2026}

\begin{document}

\begin{titlepage}
  \centering

  {\large \instituicao\par}
  \vspace{0.5cm}
  {\large \cursoAluno\par}
  \vspace{0.5cm}
  {\large \disciplina\par}

  \vfill

  {\Large \textbf{\tituloTrabalho}\par}
  \vspace{1.5cm}
  {\large Trabalho pr\'atico de banco de dados\par}

  \vfill

  \begin{flushleft}
    \textbf{Aluno:} \nomeAluno\\
    \textbf{Matr\'icula:} \matriculaAluno\\
    \textbf{Professor:} \professor\\
    \textbf{Local do minimundo:} \localMinimundo\\
    \textbf{Contato do local:} \contatoLocal
  \end{flushleft}

  \vfill

  {\large \cidadeData\par}
\end{titlepage}

\section{Descri\c{c}\~ao textual do minimundo}

O Sistema da Pastoral da Criança para Gerenciamento de Doações (SPD) tem como objetivo controlar o processo de arrecadação e distribuição de itens destinados às famílias atendidas pela instituição.
Os itens são classificados por tipos de itens. Cada tipo de item é identificado por um código e possui uma descrição. Um tipo de item pode estar associado a vários itens, enquanto cada item pertence obrigatoriamente a apenas um tipo. Os itens são identificados por um código e possuem nome, estado de conservação, tamanho e quantidade disponível.
Os doadores são identificados pelo CPF e possuem primeiro nome, sobrenome, telefone e endereço. Cada doador pode realizar diversas doações, mas cada doação é realizada por apenas um doador. As doações são identificadas por um código e possuem a data em que foram realizadas. Cada doação pode ser composta por vários itens, e um mesmo item pode estar presente em diferentes doações. Para cada item incluído em uma doação, deve ser registrada a quantidade correspondente.
As campanhas de arrecadação são identificadas por um código e possuem nome, descrição, data de início e data de término. Cada campanha pode arrecadar diversos itens, e um mesmo item pode ser arrecadado em diferentes campanhas.
Os voluntários são identificados pelo CPF e possuem nome, telefone e função desempenhada na Pastoral da Criança. Cada voluntário pode participar da organização de diversas campanhas, e cada campanha pode ser organizada por vários voluntários. Além disso, os voluntários podem participar de diversas distribuições de doações.
As famílias atendidas são identificadas por um código e possuem endereço, telefone e quantidade de membros. Cada família é composta por uma ou mais pessoas. As pessoas são identificadas pelo CPF e possuem nome, sexo e data de nascimento. Cada pessoa pertence a uma única família.
As distribuições são registradas por meio de uma data de entrega. Cada família pode receber diversas distribuições ao longo do tempo, enquanto cada distribuição é destinada a apenas uma família. Em cada distribuição podem ser entregues vários itens, e um mesmo item pode ser distribuído em diferentes distribuições. Para cada item entregue em uma distribuição, deve ser registrada a quantidade correspondente.
A quantidade disponível de cada item é calculada com base na diferença entre a quantidade total recebida por meio das doações e a quantidade total distribuída às famílias atendidas pela Pastoral da Criança. Dessa forma, o sistema permite controlar o estoque disponível e acompanhar todo o fluxo de entrada e saída dos itens arrecadados.


\section{Modelo Conceitual: Diagrama Entidade-Relacionamento em Nota\c{c}\~ao Peter Chen}

\inserirDiagrama[1.0\textwidth]{DiagramaDERTP.drawio.png}{Modelo conceitual em nota\c{c}\~ao Peter Chen}{fig:modelo-conceitual}

\section{Modelo L\'ogico: Diagrama de Esquema}

\inserirDiagrama[1.2\textwidth]{ModeloLogico.drawio.png}{Modelo l\'ogico do banco de dados}{fig:modelo-logico}

\section{Modelo de Implementa\c{c}\~ao: Diagrama de Implementa\c{c}\~ao em Nota\c{c}\~ao P\'e-de-Galinha}

\inserirDiagrama[0.95\textwidth]{PeDeGalinhaImagem.png}{Modelo de implementa\c{c}\~ao em nota\c{c}\~ao p\'e-de-galinha}{fig:modelo-implementacao}

\section{Script SQL para cria\c{c}\~ao do banco de dados}

\IfFileExists{scriptCriacao.txt}{
  \lstinputlisting[style=sqlstyle]{scriptCriacao.txt}
}{
  \begin{center}
    \fbox{
      \begin{minipage}{0.9\textwidth}
        Envie o arquivo \texttt{scriptCriacao.txt} para o Overleaf para que o
        script de cria\c{c}\~ao apare\c{c}a nesta se\c{c}\~ao.
      \end{minipage}
    }
  \end{center}
}

\section{Script SQL para popula\c{c}\~ao do banco de dados}

\IfFileExists{scriptPopular.txt}{
  \lstinputlisting[style=sqlstyle]{scriptPopular.txt}
}{
  \begin{center}
    \fbox{
      \begin{minipage}{0.9\textwidth}
        Envie o arquivo \texttt{scriptPopular.txt} para o Overleaf para que o
        script de popula\c{c}\~ao apare\c{c}a nesta se\c{c}\~ao.
      \end{minipage}
    }
  \end{center}
}

\end{document}
